\documentclass[11pt]{amsart}
\usepackage{amsmath,amssymb,amsthm}
\usepackage{graphicx}
\usepackage{tikz-cd}
\usepackage{xcolor}
\usepackage{flows}
\usepackage{loom}

\theoremstyle{plain}
\newtheorem{theorem}{Theorem}[section]
\newtheorem{lemma}[theorem]{Lemma}
\newtheorem{proposition}[theorem]{Proposition}
\newtheorem{corollary}[theorem]{Corollary}
\theoremstyle{definition}
\newtheorem{definition}[theorem]{Definition}
\newtheorem{convention}[theorem]{Convention}
\newtheorem{example}[theorem]{Example}
\theoremstyle{remark}
\newtheorem{remark}[theorem]{Remark}

\newenvironment{caveat}{\par\medskip\noindent\textbf{Caveat.}\begingroup\itshape}{\endgroup\par\medskip}
% !LOOM environment: caveat = Caveat, remark

\title{Balanced flows on finite quivers}
\author{The loom showcase}

\begin{document}
\maketitle

\begin{abstract}
We compute the rank of the lattice of balanced weightings of a finite quiver, show that it is a saturated sublattice of the weighting lattice, and count its bounded members. Nothing here is true by accident: the quiver is a showcase.
\end{abstract}

% !LOOM child: sh-0100

% !LOOM child: sh-0200

% !LOOM child: sh-0300

% !LOOM child: sh-0400

\end{document}
